\chapter{Design Decision Record}
\label{app:decisions}

This appendix records the principal design decisions of
Chapter~\ref{ch:design} as architecture decision records. Each record gives the
decision, the chosen option and its rationale, the rejected alternative, and the
consequences and evidence. The decisions are the ones a reviewer would most
challenge, and each is tied to a requirement or a contribution. Section
\ref{sec:design-decisions} names them in one sentence each.

\begingroup\setlength{\tabcolsep}{4pt}
\begin{longtable}{@{}p{0.14\textwidth}p{0.29\textwidth}p{0.21\textwidth}p{0.25\textwidth}@{}}
\caption{Architecture decision records.}\label{tab:adr}\\
\toprule
\textbf{Decision} & \textbf{Chosen option and rationale} & \textbf{Rejected alternative} & \textbf{Consequences and evidence} \\
\midrule
\endfirsthead
\toprule
\textbf{Decision} & \textbf{Chosen option and rationale} & \textbf{Rejected alternative} & \textbf{Consequences and evidence} \\
\midrule
\endhead
\bottomrule
\endfoot
ADR-01 & Correctness by certificate, not raw equality with the offline output (R1). & Comparing the candidate's tree to one offline reference tree. & Accepts any valid alternative and rejects any invalid one; generalises across all six problems. The offline oracle is kept as the efficiency baseline. Evidence: \texttt{check\_invariant} per problem; zero false negatives in the self-test. \\
\addlinespace[2pt]
ADR-02 & Isomorphism classes (nauty) for the directed universe, every labelled graph for the undirected one. & Every labelled digraph, or a random sample. & Finite, meaningful, reproducible bound, label-independent where it matters; relabellings add label-handling coverage. Evidence: \texttt{graph\_generator.py}; the universe sizes in Section~\ref{sec:universe}. \\
\addlinespace[2pt]
ADR-03 & Two-tier grade: search bound during the run, full bound on the best candidates (R3). & One bound for everything. & Seconds per candidate during the search, minutes on the finalists; closes the gap where a search-grade pass fails the full bound. Evidence: the no-reflection BFS seed that passed $n \le 4$ and failed one case at $n = 5$. \\
\addlinespace[2pt]
ADR-04 & Verifier workers as processes, not threads. & A thread pool over the chunks. & The per-chunk work is CPU-bound Python; the interpreter lock would serialise threads, so processes give real parallelism. Evidence: \texttt{ProcessPoolExecutor} in \texttt{verifier.py}; parallel equals serial in the merge test. \\
\addlinespace[2pt]
ADR-05 & Namespace jail for the sandbox, not a plain subprocess or a container runtime. & A restricted subprocess environment; a full container. & A subprocess env is not a boundary; a container adds a runtime dependency. Unprivileged user namespaces give network, PID, mount, and IPC isolation with no root and no extra packages. Evidence: \texttt{sandbox.py}; the negative tests in Chapter~\ref{ch:testing}. \\
\addlinespace[2pt]
ADR-06 & Event-sourced files, not a database. & A relational or key-value store. & Single writer, append-only history, human-inspectable artefacts, replayable, no server or migrations. Evidence: \texttt{events.jsonl}, \texttt{checkpoint.py}; resume and replay tests. \\
\addlinespace[2pt]
ADR-07 & Built-in agents plus an external coding-agent harness (R4). & Built-in agents only, or harness only. & The harness exposes the same verifier to an external agent as a baseline; the built-in loop is the registered contribution. Both share one \texttt{Backend} interface. Evidence: \texttt{agents/harness.py}; the harness arm in the evaluation. \\
\addlinespace[2pt]
ADR-08 & A pure reducer for the UI, shared by live and replay. & Separate live and replay rendering paths. & Replay is the same code on fewer events, so it cannot disagree with the live view; reconnects neither lose nor repeat events. Evidence: \texttt{ui/src/state.ts}; the Vitest reducer tests. \\
\addlinespace[2pt]
ADR-09 & Plateau-triggered specialist agents, with state derived from the event log. & Fixed agent roles only; or specialists with their own persistent state. & Dynamic generation is the registered C3 mechanism; deriving state from the log makes it resume-safe and auditable. Evidence: \texttt{specialists.py}; the spawn-and-retire test. \\
\bottomrule
\end{longtable}
\endgroup
